%%%%%%%%%%%%%%%%%%%%%%%%%%%%%%%%%%%%%%%%%%%%%%%%%%%%%%%%%%%%%%%%%%%%%%%%%%%%%%%
% 1. CONFIGURACIÓN DEL DOCUMENTO, IDIOMA Y CODIFICACIÓN
%%%%%%%%%%%%%%%%%%%%%%%%%%%%%%%%%%%%%%%%%%%%%%%%%%%%%%%%%%%%%%%%%%%%%%%%%%%%%%%

% --- Codificación y Geometría ---
% \usepackage[utf8]{inputenc} % (Recomendado si usas pdfLaTeX)
% \usepackage[T1]{fontenc}    % (Recomendado para fuentes y guiones)

\usepackage[spanish,es-tabla]{babel} % Soporte para idioma español, separación silábica y nombres
\usepackage{csquotes}                % Soporte avanzado para comillas y citas (funciona con babel)
\usepackage[a4paper, margin=2cm, marginparwidth=3cm]{geometry} % Configura márgenes y tamaño de página

%%%%%%%%%%%%%%%%%%%%%%%%%%%%%%%%%%%%%%%%%%%%%%%%%%%%%%%%%%%%%%%%%%%%%%%%%%%%%%%
% 2. PAQUETES DE MATEMÁTICA Y FÍSICA
%%%%%%%%%%%%%%%%%%%%%%%%%%%%%%%%%%%%%%%%%%%%%%%%%%%%%%%%%%%%%%%%%%%%%%%%%%%%%%%

\usepackage{amsmath}   % Comandos avanzados para entornos matemáticos
\usepackage{amssymb}   % Símbolos matemáticos extendidos (como \mathbb)
\usepackage{bbold}     % Fuentes "Blackboard" (para números 0, 1)
\usepackage{bigints}   % Símbolos de integrales más grandes
\usepackage{physics}   % Notación común en física (derivadas, \bra, \ket, etc.)
\usepackage{cancel}    % Para tachar expresiones matemáticas
\usepackage{bm}        % Para usar negrita en modo matemática

%%%%%%%%%%%%%%%%%%%%%%%%%%%%%%%%%%%%%%%%%%%%%%%%%%%%%%%%%%%%%%%%%%%%%%%%%%%%%%%
% 3. GRÁFICOS, FIGURAS, COLORES Y DIBUJO (TIKZ)
%%%%%%%%%%%%%%%%%%%%%%%%%%%%%%%%%%%%%%%%%%%%%%%%%%%%%%%%%%%%%%%%%%%%%%%%%%%%%%%

\usepackage[table,xcdraw,x11names]{xcolor}      % Gestión avanzada de colores (usado por TikZ, listings, etc.)
\usepackage{graphicx}                           % Para incluir imágenes (.png, .jpg, .pdf)
\usepackage[export]{adjustbox}                  % Herramientas extra para imágenes (recortar, enmarcar)
\usepackage{float}                              % Mejor control sobre la posición de flotantes [H]
\usepackage[font=small, labelfont=bf]{caption}  % Personalización de los pies de foto/tabla
\usepackage{subcaption}                         % Para crear subfiguras y subtables
\usepackage{tikz}                               % Potente sistema para crear gráficos vectoriales
\usepackage{circuitikz}                         % Extensión de TikZ para dibujar circuitos eléctricos

% --- Librerías y definiciones de TikZ ---
\usetikzlibrary{patterns, decorations.markings, arrows.meta}
\usetikzlibrary{decorations.pathreplacing}
\usetikzlibrary{decorations.pathmorphing}
\definecolor{ForestGreen}{rgb}{0.13, 0.55, 0.13}    % Define un color personalizado
\definecolor{Yellow}{HTML}{91760D}                  % Define un color personalizado


%%%%%%%%%%%%%%%%%%%%%%%%%%%%%%%%%%%%%%%%%%%%%%%%%%%%%%%%%%%%%%%%%%%%%%%%%%%%%%%
% 4. FORMATO DE PÁGINA, SECCIONES Y LISTAS
%%%%%%%%%%%%%%%%%%%%%%%%%%%%%%%%%%%%%%%%%%%%%%%%%%%%%%%%%%%%%%%%%%%%%%%%%%%%%%%

\usepackage{titlesec}     % Personalización avanzada de títulos de sección
\usepackage{multicol}     % Para crear texto en múltiples columnas
\usepackage{enumitem}     % Personalización avanzada de listas (enumerate, itemize)
\usepackage{wrapfig}      % Para que el texto fluya alrededor de una figura
\usepackage{booktabs}     % Crea tablas con aspecto profesional (mejores líneas)
\usepackage{setspace}     % Permite cambiar el interlineado (ej. \onehalfspacing)
\usepackage{yfonts}       % Permite usar fuentes góticas (\textfrak)
\usepackage{fontawesome}  % Íconos de Font Awesome

%%%%%%%%%%%%%%%%%%%%%%%%%%%%%%%%%%%%%%%%%%%%%%%%%%%%%%%%%%%%%%%%%%%%%%%%%%%%%%%
% 5. UTILIDADES Y ENTORNOS ESPECIALES
%%%%%%%%%%%%%%%%%%%%%%%%%%%%%%%%%%%%%%%%%%%%%%%%%%%%%%%%%%%%%%%%%%%%%%%%%%%%%%%

\usepackage{todonotes}          % Añade notas (TODOs) en el margen o en línea
\usepackage[most]{tcolorbox}    % Crea cajas de texto coloreadas y personalizables
\usepackage{listings}           % Para insertar bloques de código fuente (Syntax Highlighting)
\usepackage{pdfpages}           % Para incluir páginas de otros documentos PDF
\usepackage{lipsum}             % Genera texto de relleno (Lorem Ipsum)
\usepackage{bm}

%%%%%%%%%%%%%%%%%%%%%%%%%%%%%%%%%%%%%%%%%%%%%%%%%%%%%%%%%%%%%%%%%%%%%%%%%%%%%%%
% 6. ENCABEZADOS, PIES DE PÁGINA E ÍNDICES (TOC)
%%%%%%%%%%%%%%%%%%%%%%%%%%%%%%%%%%%%%%%%%%%%%%%%%%%%%%%%%%%%%%%%%%%%%%%%%%%%%%%

\usepackage{tocloft}      % Personalización de la Tabla de Contenidos (ToC)
\usepackage{fancyhdr}     % Para personalizar encabezados y pies de página
\usepackage{lastpage}     % Provee una referencia a la última página (para "Pág X de Y")

%%%%%%%%%%%%%%%%%%%%%%%%%%%%%%%%%%%%%%%%%%%%%%%%%%%%%%%%%%%%%%%%%%%%%%%%%%%%%%%
% 7. HIPERVÍNCULOS (IMPORTANTE: Cargar casi al final)
%%%%%%%%%%%%%%%%%%%%%%%%%%%%%%%%%%%%%%%%%%%%%%%%%%%%%%%%%%%%%%%%%%%%%%%%%%%%%%%

% Generalmente es mejor cargar hyperref al final, ya que modifica
% muchos comandos de LaTeX.
\usepackage[colorlinks=true, allcolors=black]{hyperref} 

%%%%%%%%%%%%%%%%%%%%%%%%%%%%%%%%%%%%%%%%%%%%%%%%%%%%%%%%%%%%%%%%%%%%%%%%%%%%%%%
% 8. CONFIGURACIONES Y COMANDOS PERSONALIZADOS
% (Deben ir después de cargar sus paquetes correspondientes)
%%%%%%%%%%%%%%%%%%%%%%%%%%%%%%%%%%%%%%%%%%%%%%%%%%%%%%%%%%%%%%%%%%%%%%%%%%%%%%%

% --- Config. Mates (amsmath) ---
% Numera las ecuaciones como (sección.numero)
\numberwithin{equation}{section} 

% --- Config. Secciones y TOC (titlesec, babel) ---
% Formato para párrafos como si fueran subtítulos
\titleformat{\paragraph}[block]{\normalfont\normalsize\bfseries}{\theparagraph}{1em}{}
\setcounter{secnumdepth}{4} % Profundidad de numeración de secciones
\setcounter{tocdepth}{4}    % Profundidad del índice (TOC)

% --- Config. Encabezados (fancyhdr) ---
\pagestyle{fancy}                                   % Activa el estilo personalizado
\fancyhf{}                                          % Limpia todos los campos de encabezado y pie
\setlength{\headheight}{16pt}                       % Aumenta el espacio para el encabezado
\fancyhead[R]{\rightmark}                           % Pone el nombre de la sección/subsección a la derecha
\fancyhead[L]{\hyperlink{toc}{Índice}}              % Enlace al TOC a la izquierda (requiere hyperref)
\cfoot{\thepage \hspace{1pt} de \pageref{LastPage}} % Pie de página centrado "Página X de Y"

% Asegura que las subsecciones actualicen \rightmark (para el encabezado)
\addto\captionsspanish{\renewcommand{\subsectionmark}[1]{\markright{#1}}}

% --- Config. Layout (multicol) ---
\setlength\columnsep{18pt}          % Separación entre columnas

% --- Config. Utilidades (todonotes) ---
\setuptodonotes{inline, size=\footnotesize, color=blue!20, bordercolor=none}

% --- Config. Código (listings) ---
\lstdefinelanguage{custombash}{
    language=bash,
    morekeywords={conda, activate, mkdir, cd, cmake, make, git, sudo},
    sensitive=true
}

\lstdefinestyle{bashStyle}{
    language=custombash,
    basicstyle=\ttfamily\small,
    keywordstyle=\color{teal}\bfseries,
    commentstyle=\color{gray}\itshape,
    stringstyle=\color{orange},
    identifierstyle=\color{black},
    numbers=left,
    numberstyle=\tiny\color{gray},
    stepnumber=1,
    numbersep=10pt,
    backgroundcolor=\color{white},
    showspaces=false,
    showstringspaces=false,
    breaklines=true,
    frame=single,
    captionpos=b
}

\lstdefinestyle{pythonStyle}{
    language=Python,
    basicstyle=\ttfamily\small,
    keywordstyle=\color{blue}\bfseries,
    commentstyle=\color{gray}\itshape,
    stringstyle=\color{orange},
    identifierstyle=\color{black},
    numbers=left,
    numberstyle=\tiny\color{gray},
    stepnumber=1,
    numbersep=10pt,
    backgroundcolor=\color{white},
    showspaces=false,
    showstringspaces=false,
    breaklines=true,
    frame=single,
    captionpos=b
}

\lstdefinestyle{cppStyle}{
    language=C++,
    basicstyle=\ttfamily\small,
    keywordstyle=\color{blue}\bfseries,
    commentstyle=\color{gray}\itshape,
    stringstyle=\color{red},
    numbers=left,
    numberstyle=\tiny\color{gray},
    stepnumber=1,
    numbersep=10pt,
    backgroundcolor=\color{white},
    showspaces=false,
    showstringspaces=false,
    breaklines=true,
    frame=single,
    captionpos=b
}

% --- Comandos Personalizados ---

% Comando para texto en un círculo (requiere TikZ)
\newcommand*\circled[1]{\tikz[baseline=(char.base)]{\node[shape=circle,draw,inner sep=1.5pt] (char) {#1};}}

% Comando para tachar con color (requiere cancel y xcolor)
\newcommand{\colorcancel}[2]{\textcolor{#1}{\cancel{\textcolor{black}{#2}}}}

% --- Entornos Personalizados (tcolorbox) ---
\newtcolorbox{infobox}[2][blue]{
    colback=#1!5!white,
    colframe=#1!75!black,
    title={#2},
    fonttitle=\bfseries,
    coltitle=black,
    boxrule=0.8pt,
    arc=4pt,
    left=6pt,
    right=6pt,
    top=6pt,
    bottom=6pt,
    breakable
}

%%%%%%%%%%%%%%%%%%%%%%%%%%%%%%%%%%%%%%%%%%%%%%%%%%%%%%%%%%%%%%%%%%%%%%%%%%%%%%%